% !TeX program = xelatex
% Απαιτεί XeLaTeX (fontspec/polyglossia) -- ΔΕΝ μεταγλωττίζεται με pdflatex.
\documentclass[11pt,a4paper]{article}

\usepackage{fontspec}
\usepackage{polyglossia}
\setmainlanguage{greek}
\setotherlanguage{english}
\setmainfont{Linux Libertine O}
\setsansfont{Linux Biolinum O}
\setmonofont{DejaVu Sans Mono}[Scale=0.85]

\usepackage[a4paper,margin=2.5cm]{geometry}
\usepackage{parskip}
\usepackage{amsmath,amssymb}
\usepackage{booktabs}
\usepackage{array}
\usepackage{float}
\usepackage{enumitem}
\usepackage{caption}
\usepackage{xcolor}

\newcommand{\file}[1]{\texttt{#1}}
\let\oldunderscore\_
\renewcommand{\_}{\oldunderscore\allowbreak}
\setlength{\emergencystretch}{3em}
\setlist[itemize]{topsep=3pt,itemsep=2pt}

\title{\textbf{Αναφορά Κατάστασης Διπλωματικής}\\[8pt]
\large\normalfont Κυρτή Βελτιστοποίηση σε Πραγματικό Χρόνο Βασισμένη στην
Παραμετροποίηση Μείγματος Ειδικών\\ για την Πρόβλεψη Αποτελεσμάτων Αθλητικών Αγώνων}
\author{[Στέφανος Μουσάδης]}
\date{\today}

\begin{document}
\maketitle

% =====================================================================
\section{Θέμα και δεδομένα}
% =====================================================================

Η ιδέα είναι να συγκρίνουμε bookmakers μεταξύ τους και ταυτόχρονα να χτίσουμε ένα μοντέλο
πρόβλεψης πάνω τους, με online κυρτή βελτιστοποίηση, σαν μείγμα από experts. Δουλέψαμε σε
ποδόσφαιρο, με δεδομένα από το football-data.co.uk.

Έχουμε 22 πρωταθλήματα, 10 σεζόν (2016/17--2025/26) και 76.583 αγώνες όπου υπάρχει
τουλάχιστον μία χρήσιμη απόδοση. Μετά τον καθαρισμό μείναμε με 24 bookmakers: πετάξαμε τις
συγκεντρωτικές στήλες, που δεν είναι ανεξάρτητοι experts, ενώσαμε το VC Bet με το
BetVictor αφού πρόκειται για τον ίδιο bookmaker πριν και μετά το rebrand, και βγάλαμε
εντελώς την Interwetten, που σταματάει να δίνει αποδόσεις τον Ιανουάριο του 2024 αλλά στη
δεκαετία περνάει κάθε κατώφλι κάλυψης, οπότε χάλαγε σιωπηλά κάθε σύγκριση σε πρόσφατο
παράθυρο. Το βασικό πρόβλημα των δεδομένων είναι η κάλυψη: άλλος καλύπτει το 99,8\% των
αγώνων και άλλος το 6,9\%. Γι' αυτό δουλεύουμε με Sleeping Experts.

% =====================================================================
\section{Τι έγινε}
% =====================================================================

\begin{center}
\begin{minipage}{0.92\textwidth}
\ttfamily\small
220 raw CSV $\;\to\;$ de-vig $\times$ 2 (απλή / Shin) $\;\to\;$
Hedge / OGD / FTRL / Uniform\\
\hspace*{1em}$\to$ βαθμονόμηση (online temperature scaling)
$\;\to\;$ κατάταξη, bootstrap tests\\
\hspace*{1em}$\to$ Kelly με online πολλαπλασιαστή
\end{minipage}
\end{center}

\begin{itemize}
  \item Το de-vig το κάναμε με δύο μεθόδους, απλή αναλογική κανονικοποίηση και Shin, και
        κρατήσαμε δύο ξεχωριστά datasets. Θέλαμε να ξέρουμε αν κάποιο συμπέρασμα κρέμεται
        από αυτή την επιλογή.
  \item Τρέξαμε Hedge, OGD και FTRL, συν uniform baseline, σε τρία panel: μόνο opening, μόνο
        closing, όλα μαζί. Επί δύο κανονικοποιήσεις.
  \item Το μείγμα δεν εκπαιδεύεται με τη συνήθη έννοια: είναι online, κάθε πρόβλεψη
        χρησιμοποιεί μόνο τους προηγούμενους γύρους. Μετά μπαίνει βαθμονόμηση με online
        temperature scaling, δεύτερο στάδιο OCO. Ο ρυθμός της επιλέγεται σε χρονολογικό
        πρόθεμα, στο πρώτο 75\% των αγώνων κάθε σειράς, και μετά γίνεται ένα αιτιώδες
        πέρασμα σε όλη την ιστορία με αυτόν σταθερό.
  \item Κάθε σύγκριση περνάει από moving block bootstrap. Οι διαδοχικές απώλειες δεν είναι
        ανεξάρτητες, οπότε ένα σκέτο t-test θα υποεκτιμούσε τη διασπορά.
  \item Ο οικονομικός έλεγχος είναι Kelly με leave-one-bookmaker-out: ο στόχος βγαίνει από
        το μείγμα πριν την εκπαίδευση, και το panel ταιριάζει στη φάση του. Το μέγεθος του
        στοιχήματος δεν το καρφώνουμε πλέον σε σταθερό κλάσμα, το μαθαίνουμε κι αυτό online
        --- τρίτο στάδιο OCO πάνω στα δύο προηγούμενα. Το σταθερό $\tfrac14$ κρατιέται ως
        σημείο αναφοράς.
  \item Τέλος συντονίσαμε το βήμα και των τριών αλγορίθμων, στο ίδιο πλέγμα, με επιλογή στο
        πρώτο 75\% και μέτρηση στο υπόλοιπο. Και τα τρία ήθελαν αλλαγή: το OGD βήμα 100 φορές
        μεγαλύτερο από το θεωρητικό, το Hedge τέσσερις φορές μικρότερο και το FTRL είκοσι.
\end{itemize}

% =====================================================================
\section{Αποτελέσματα}
% =====================================================================

Στην κατάταξη μπαίνουν μόνο bookmakers με κάλυψη $\geq$70\%. Τους υπόλοιπους θα τους
κρίναμε σε λίγες χιλιάδες αγώνες, συχνά μιας μόνο σεζόν, και η απόδοσή τους θα μπερδευόταν
με τη δυσκολία εκείνης της σεζόν.

\begin{table}[H]
\centering
\small
\begin{tabular}{@{}lllrrr@{}}
\toprule
& & \textbf{φάση} & \textbf{log-loss} & \textbf{Brier} & \textbf{accuracy} \\
& & & (raw / cal.) & (raw / cal.) & \\
\midrule
PSC (Pinnacle) & bookmaker & closing & 0,99714 / 0,99685 & 0,59592 / 0,59577 & 50,61\% \\
OGD            & αλγόριθμος & ---    & 0,99739 / 0,99695 & 0,59609 / 0,59587 & 50,70\% \\
Hedge          & αλγόριθμος & ---    & 0,99830 / 0,99768 & 0,59666 / 0,59634 & 50,65\% \\
FTRL           & αλγόριθμος & ---    & 0,99831 / 0,99769 & 0,59667 / 0,59634 & 50,65\% \\
Uniform        & baseline   & ---    & 0,99873 / 0,99808 & 0,59695 / 0,59661 & 50,60\% \\
VC\_BV         & bookmaker & opening & 1,00075 / 1,00026 & 0,59833 / 0,59806 & 50,36\% \\
WH             & bookmaker & opening & 1,00119 / 1,00048 & 0,59860 / 0,59818 & 50,31\% \\
PS (Pinnacle)  & bookmaker & opening & 1,00083 / 1,00055 & 0,59838 / 0,59824 & 50,31\% \\
BW             & bookmaker & opening & 1,00140 / 1,00065 & 0,59873 / 0,59833 & 50,36\% \\
B365           & bookmaker & opening & 1,00113 / 1,00073 & 0,59852 / 0,59832 & 50,38\% \\
\bottomrule
\end{tabular}
\caption{Κατάταξη Tier-A στο πλήρες panel, 10 σεζόν, κατά βαθμονομημένο log-loss. Το Brier
βγάζει την ίδια σειρά με μόνη εξαίρεση τις δύο τελευταίες θέσεις, όπου B365 και BW
αντιμετατίθενται. Η accuracy βγάζει άλλη σειρά, αλλά καμία διαφορά της δεν είναι σημαντική.}
\end{table}

Προσοχή στη στήλη «φάση»: το κατώφλι του 70\% αφήνει μέσα μόνο έναν closing bookmaker, τον
PSC, γιατί οι υπόλοιπες closing στήλες αρχίζουν γύρω στο 2019/20 (ο B365C χάνει στο τσακ με
69,6\%), οπότε εδώ «ξεπερνάμε όλους εκτός του PSC» σημαίνει ότι ξεπερνάμε πέντε τιμές
ανοίγματος. Η σοβαρή σύγκριση, απέναντι σε πραγματικές τιμές κλεισίματος, γίνεται στην
τελευταία ενότητα.

\begin{itemize}
  \item Και οι τρεις αλγόριθμοι ξεπερνούν σημαντικά το uniform baseline: OGD $-0{,}00134$,
        Hedge $-0{,}00044$, FTRL $-0{,}00042$, όλα $p<0{,}001$.
  \item Το OGD ξεπερνά σημαντικά Hedge, FTRL \emph{και} B365C ($-0{,}00077$).
  \item Η βαθμονόμηση βοηθάει σημαντικά παντού, $+0{,}00045$ έως $+0{,}00066$.
  \item Ο PSC μένει μπροστά μας εδώ, και η διαφορά είναι σημαντική: $+0{,}00044$
        βαθμονομημένο, στους 71.818 αγώνες που καλύπτει.
\end{itemize}

\subsection*{Ποιο panel}

\begin{table}[H]
\centering
\small
\begin{tabular}{@{}lrrrrrr@{}}
\toprule
& \multicolumn{2}{c}{\textbf{opening} (12)} & \multicolumn{2}{c}{\textbf{όλα μαζί} (24)} & \multicolumn{2}{c}{\textbf{closing} (12)} \\
\cmidrule(lr){2-3}\cmidrule(lr){4-5}\cmidrule(lr){6-7}
& απλή & Shin & απλή & Shin & απλή & Shin \\
\midrule
OGD     & 1,00068 & 1,00021 & 0,99739 & 0,99707 & 0,99714 & 0,99685 \\
Hedge   & 1,00079 & 1,00029 & 0,99830 & 0,99774 & 0,99758 & 0,99713 \\
FTRL    & 1,00079 & 1,00029 & 0,99831 & 0,99774 & 0,99758 & 0,99713 \\
Uniform & 1,00071 & 1,00018 & 0,99873 & 0,99816 & 0,99743 & 0,99701 \\
\bottomrule
\end{tabular}
\caption{Log-loss ανά panel και κανονικοποίηση, 10 σεζόν.}
\end{table}

Η σειρά closing $<$ όλα μαζί $<$ opening κρατάει σε κάθε κελί, και στις 24 συγκρίσεις που
τρέξαμε η διαφορά βγαίνει σημαντική. Ότι κερδίζουν τα closing το περιμέναμε. Αυτό που δεν
περιμέναμε είναι ότι το closing-only κερδίζει και το \emph{πλήρες} panel: οι opening
bookmakers δεν προσφέρουν διαφοροποίηση, απλώς βαραίνουν το μείγμα. Το Shin ανεβάζει όλα τα
κελιά κατά 0,0003 ως 0,0006 χωρίς να αλλάζει καμία σύγκριση.

Κοιτάξαμε και τι κάνει κάθε αλγόριθμος απέναντι στον \emph{απλό μέσο όρο}, μέσα στο ίδιο
panel:

\begin{itemize}
  \item Στο πλήρες panel κερδίζουν σημαντικά και οι τρεις, με το OGD να έχει περίπου
        τριπλάσιο περιθώριο από τους άλλους δύο.
  \item Στα περιορισμένα panel τα Hedge και FTRL \emph{χάνουν σημαντικά} από τον μέσο όρο,
        8 στα 8 τεστ. Το OGD κερδίζει στο closing και ισοπαλεί στο opening.
\end{itemize}

Το βλέπουμε έτσι: το πλήρες panel έχει 12 καλούς closing και 12 σταθερά χειρότερους opening,
δηλαδή μια ποιοτική διαφορά που αξίζει να τη μάθεις. Σε ομοιογενές panel αυτή η διαφορά
λείπει, και η προσαρμογή κοστίζει περισσότερο απ' όσο αποδίδει.

\subsection*{Η διαμόρφωση που κρατάμε}

Είχαμε τρεις βελτιώσεις μετρημένες χωριστά (το διορθωμένο βήμα, το closing-only panel, το
Shin) και τις στοιβάξαμε. Η παγίδα εδώ είναι ότι τα panel καλύπτουν διαφορετικούς αγώνες,
και ένα log-loss πάνω σε άλλο σύνολο αγώνων δεν συγκρίνεται με τίποτα. Βαθμολογήσαμε λοιπόν
κάθε διαμόρφωση στους ίδιους 71.818 αγώνες, την τομή όλων με τον PSC. Διαλέξαμε στο πρώτο
75\% και κρατήσαμε την ουρά στην άκρη, γιατί το να διαλέγεις κοιτώντας αυτά τα νούμερα είναι
κι αυτό επιλογή.

\begin{table}[H]
\centering
\begin{tabular}{@{}lrr@{}}
\toprule
\textbf{Διαμόρφωση} & \textbf{πλήρες δείγμα} & \textbf{held-out} \\
\midrule
PSC μόνος του                             & 0,99685 & 0,99502 \\
closing-only + Shin \emph{(αυτή κρατάμε)} & 0,99691 & 0,99505 \\
closing-only, απλή                        & 0,99692 & 0,99503 \\
όλα μαζί, απλή                            & 0,99729 & 0,99504 \\
όλα μαζί, Shin                            & 0,99731 & 0,99506 \\
\bottomrule
\end{tabular}
\caption{Βαθμονομημένο log-loss, ίδιοι 71.818 αγώνες για κάθε γραμμή.}
\end{table}

\begin{itemize}
  \item Όλο το κέρδος έρχεται από το panel. Το closing-only $+$ Shin κερδίζει το πλήρες κατά
        $0{,}00038$ ($p<0{,}001$), ενώ το Shin σκέτο πάνω στο πλήρες panel δεν αλλάζει τίποτα
        μετρήσιμο ($p=0{,}253$).
  \item Εδώ ισοπαλούμε με τον PSC: $+0{,}00006$, $p=0{,}227$. Στο πλήρες panel χάναμε
        σημαντικά.
  \item Στην held-out ουρά και οι πέντε γραμμές μαζεύονται σε λιγότερο από $0{,}00005$. Η
        ουρά επιβεβαιώνει την ισοπαλία, δεν βγάζει νικητή, και δεν θα ήταν σωστό να την
        παρουσιάσουμε σαν νίκη.
\end{itemize}

\subsection*{Χωρίς την Pinnacle}

Το προφανές αντεπιχείρημα σε όλα αυτά είναι ότι το μείγμα απλώς αναπαράγει την Pinnacle, που
είναι ούτως ή άλλως μέσα του. Βγάλαμε λοιπόν και τις δύο στήλες της, PS και PSC (είναι η ίδια
εταιρεία, opening και closing), και ξαναεκπαιδεύσαμε από την αρχή σε τρία panel. Η απάντηση
δεν είναι μία.

\begin{table}[H]
\centering
\small
\begin{tabular}{@{}llrll@{}}
\toprule
\textbf{Panel} & \textbf{Θέση OGD} & \textbf{log-loss} & \textbf{Πρώτος} & \textbf{Ζευγαρωτά} \\
\midrule
πλήρες (22 experts)      & \#3 από 8 & 0,99769 & FTRL 0,99750    & 4 νίκες, 0 ήττες \\
closing-only (11)        & \#5 από 8 & 0,99866 & Uniform 0,99857 & 3 νίκες, 1 ισοπαλία \\
opening-only (11)        & \#5 από 8 & 0,99988 & VC\_BV 0,99969  & 3 νίκες, \textbf{1 ήττα} \\
\bottomrule
\end{tabular}
\caption{Κάθε panel ξαναεκπαιδεύεται από την αρχή και συγκρίνεται με τους bookmakers της
δικής του φάσης, όλοι στους ίδιους αγώνες (58.567 για πλήρες και opening, 35.852 για
closing).}
\end{table}

Στο \textbf{πλήρες} panel αντέχει άνετα. Και οι τέσσερις μέθοδοι βγαίνουν πάνω από όλους τους
bookmakers, οι τρεις αλγόριθμοι απέχουν μόλις 0,0002 μεταξύ τους χωρίς καμία σημαντική
διαφορά μεταξύ τους, και το OGD κερδίζει σημαντικά και τους τέσσερις ($-0{,}0019$ έως
$-0{,}0027$, $p<0{,}001$ παντού). Με έναν αστερίσκο: όλοι τους είναι opening.

Στα δύο ομοιογενή panel τα πράγματα αλλάζουν, και αλλάζουν με τον ίδιο ακριβώς τρόπο. Στο
closing το OGD κερδίζει τρεις (WHC $-0{,}00093$, BWC $-0{,}00049$, B365C $-0{,}00021$) και
ισοπαλεί με τον VC\_BVC. Στο opening κερδίζει πάλι τρεις, αλλά \textbf{χάνει σημαντικά από
τον VC\_BV} ($+0{,}00025$, $p=0{,}009$) --- η μόνη φορά σε όλη την εργασία που μας ξεπερνά
bookmaker εκτός της Pinnacle. Και στα δύο ο απλός μέσος όρος περνάει μπροστά του.

Το μοτίβο είναι αρκετά καθαρό ώστε να μην το αποδώσουμε στην τύχη: μόλις φύγει η Pinnacle,
οι υπόλοιποι της ίδιας φάσης μοιάζουν τόσο πολύ μεταξύ τους που δεν μένει ποιοτική διαφορά
για να τη μάθει η στάθμιση, και ο σκέτος μέσος όρος αρκεί. Όταν το είχαμε δει μόνο στο
closing μπορούσε να περάσει για ιδιοτροπία εκείνου του panel. Τώρα που εμφανίζεται
πανομοιότυπα και στο opening, είναι ιδιότητα της ομοιογένειας.

Η διαμόρφωση που κρατάμε είναι λοιπόν η καλύτερη εν μέρει επειδή περιέχει τον καλύτερο
odds-setter της αγοράς. Δεν το κρύβουμε.

Ότι δεν ξεπερνάμε την Pinnacle δεν οφείλεται σε καλύτερο μοντέλο πρόβλεψης από τη δική της
μεριά, αλλά σε διαφορετικό επιχειρηματικό μοντέλο. Δουλεύει με χαμηλό περιθώριο και υψηλό
τζίρο, και δεν κόβει τους κερδισμένους παίκτες.

\subsection*{Πρόβλεψη για το 2025--26, εκπαιδευμένη ως το 2024}

Κάθε νούμερο ως εδώ βαθμολογεί το μείγμα πάνω στην ίδια ιστορία που έμαθε. Αυτό είναι θεμιτό
για OCO, αφού ο αλγόριθμος είναι αιτιώδης και η σωρευτική απώλεια είναι ακριβώς το μέγεθος
που φράσσει το regret, αλλά δεν είναι η ερώτηση που κάνει κανείς στην πράξη. Κόψαμε λοιπόν
σε πραγματική ημερομηνία: εκπαίδευση στους 65.009 αγώνες ως 31/12/2024, μέτρηση στους
11.574 του 2025--26. Αντίπαλοι όσοι δίνουν αποδόσεις \emph{εκεί}, με κατώφλι υπολογισμένο
μέσα στο test.

\begin{table}[H]
\centering
\small
\begin{tabular}{@{}rllr@{}}
\toprule
\textbf{\#} & \textbf{Σειρά} & \textbf{φάση} & \textbf{log-loss} \\
\midrule
1 & BFEC (Betfair Exchange) & closing & 1,00040 \\
2 & \textbf{OGD closing-only, online} & --- & \textbf{1,00109} \\
3 & OGD πλήρες panel, online & --- & 1,00128 \\
4 & BWC & closing & 1,00146 \\
5 & OGD closing-only, παγωμένο & --- & 1,00151 \\
6 & B365C & closing & 1,00165 \\
7 & OGD πλήρες panel, παγωμένο & --- & 1,00183 \\
8 & B365 & opening & 1,00371 \\
9 & BW & opening & 1,00416 \\
10 & BFE & opening & 1,00758 \\
\bottomrule
\end{tabular}
\caption{8.151 κοινοί αγώνες του 2025--26. «Online» σημαίνει ότι το μείγμα συνεχίζει να
μαθαίνει μέσα στο test, «παγωμένο» ότι τα βάρη κλειδώνουν στην τομή.}
\end{table}

\begin{itemize}
  \item \textbf{Κάθε τιμή ανοίγματος την κερδίζουμε σημαντικά}, σε κάθε παραλλαγή:
        $-0{,}0022$ έως $-0{,}0065$, όλα $p\leq0{,}007$.
  \item \textbf{Τις τιμές κλεισίματος δεν τις κερδίζουμε.} Το καλύτερο που πετυχαίνουμε
        είναι ισοπαλία και με τις τρεις, και μόνο με μία παραλλαγή (closing-only, online):
        B365C $p=0{,}184$, BWC $p=0{,}334$, BFEC $p=0{,}051$. Στις άλλες τρεις ο BFEC μάς
        κερδίζει σημαντικά.
  \item Το online κερδίζει σταθερά το παγωμένο ($1{,}00109$ έναντι $1{,}00151$), οπότε η
        σωστή διατύπωση δεν είναι «εκπαιδεύτηκε ως το 2024» αλλά «τρέχει συνεχώς».
\end{itemize}

\paragraph{Σημείωση για την Pinnacle.} Λείπει από τον πίνακα γιατί η κάλυψή της πέφτει στο
59\% μετά το 2025 και σταματάει να δίνει αποδόσεις στις 14/01/2026· αν την απαιτούσαμε σε
κάθε γραμμή, οι κοινοί αγώνες θα έπεφταν από 8.151 σε 2.860 και δεν θα έβγαινε τίποτα
σημαντικό πουθενά. Τη συγκρίναμε λοιπόν ένα προς ένα, στους αγώνες που καλύπτει η ίδια, και
το αποτέλεσμα είναι σταθερό και στις τέσσερις παραλλαγές: \textbf{τον PS τον κερδίζουμε
σημαντικά} ($-0{,}0022$ έως $-0{,}0023$, $p$ από $0{,}006$ ως $0{,}015$, 6.857 αγώνες) και
\textbf{με τον PSC ισοπαλούμε} ($-0{,}0001$ έως $-0{,}0002$, $p$ από $0{,}149$ ως $0{,}549$,
6.892 αγώνες). Το σημείο αναφοράς της αγοράς σήμερα δεν είναι πάντως η Pinnacle αλλά το
Betfair Exchange closing, που δεν είναι bookmaker αλλά ανταλλακτήριο, δηλαδή τιμή αγοράς
χωρίς περιθώριο διαχειριστή.

Μετρήσαμε τέλος αν το μείγμα παγώνει λόγω του βήματος $1/\sqrt{t}$: δεν παγώνει. Η μέση
κίνηση των βαρών ανά γύρο είναι $0{,}0226$, $0{,}0156$, $0{,}0138$, $0{,}0157$ στα τέσσερα
τεταρτημόρια, γιατί κυριαρχεί η προβολή στο simplex κάθε φορά που αλλάζει το σύνολο των
ξύπνιων experts, όχι το βήμα.

\subsection*{Value betting και το μέγεθος του στοιχήματος}

Το panel εκπαίδευσης το ταιριάζουμε στη φάση του στόχου: closing-only για closing στόχο,
opening-only για opening. Το δεύτερο δεν είναι επιλογή αλλά απαίτηση, γιατί καμία closing
τιμή δεν υπάρχει ακόμα τη στιγμή που βγαίνει μια τιμή ανοίγματος.

Το μέγεθος του στοιχήματος το μαθαίνουμε πλέον κι αυτό. Ως τώρα ήταν καρφωμένο στο ένα
τέταρτο του Kelly, που είναι πρακτικός κανόνας και τίποτα παραπάνω. Τώρα κρατάμε τον τύπο
του Kelly, ο οποίος ξέρει την πιθανότητα και την απόδοση του συγκεκριμένου αγώνα, και
αφήνουμε τον αλγόριθμο να μάθει μόνο τον \emph{πολλαπλασιαστή}: πόσο να τον συρρικνώσει.
Ξεκινάει από το πλήρες Kelly και τον κατεβάζει όσο τον διαψεύδουν τα αποτελέσματα.

Το είχαμε ξαναδοκιμάσει, αλλά τότε μαθαίναμε απευθείας το κλάσμα του κεφαλαίου, πετώντας
την πιθανότητα που είχαν παραγάγει τα δύο προηγούμενα στάδια και ξαναμαθαίνοντάς την από τα
αποτελέσματα. Εκείνη η εκδοχή έχανε συστηματικά: 7 ήττες στις 20 συγκρίσεις, καμία νίκη. Με
τον πολλαπλασιαστή η εικόνα αλλάζει τελείως.

\begin{table}[H]
\centering
\small
\begin{tabular}{@{}lrrrrr@{}}
\toprule
\textbf{Εναντίον} & \textbf{Στοιχ.} & \textbf{Kelly $\tfrac14$} & \textbf{OCO λ} &
\textbf{Kelly πλήρες} & \textbf{μέσο λ} \\
\midrule
WHC     &  7.216 & $84{,}3\times$ & $\mathbf{5.945\times}$ & $28.079\times$ & 0,79 \\
BWC     &  6.868 & $2{,}82\times$ & $\mathbf{5{,}27\times}$ & $9{,}65\times$ & 0,63 \\
VC\_BVC &  4.735 & $2{,}01\times$ & $\mathbf{2{,}87\times}$ & $4{,}31\times$ & 0,61 \\
B365C   &  7.361 & $1{,}26\times$ & $1{,}02\times$ & $0{,}51\times$ & 0,47 \\
PSC     & 20.198 & $0{,}38\times$ & $0{,}20\times$ & $0{,}0008\times$ & 0,29 \\
\bottomrule
\end{tabular}
\caption{Τελικό κεφάλαιο με αρχή 1, βαθμονομημένες πιθανότητες, απλή κανονικοποίηση, ίδια
στοιχήματα σε κάθε στήλη. Η τελευταία στήλη είναι ο μέσος πολλαπλασιαστής που έμαθε ο
αλγόριθμος· η στήλη Kelly $\tfrac14$ αντιστοιχεί σε σταθερό 0,25.}
\end{table}

Τα δύο σκέλη του αποτελέσματος δεν λένε το ίδιο πράγμα. Στον ζευγαρωμένο έλεγχο ανά
στοίχημα \textbf{ισοπαλούμε} με το σταθερό ένα τέταρτο: 18 ισοπαλίες, μία νίκη, μία ήττα
στα 20 κελιά. Στο κεφάλαιο όμως ο μαθαινόμενος πολλαπλασιαστής είναι μπροστά σε τρεις από
τους πέντε στόχους, και στον WHC κατά εβδομήντα φορές. Τα δύο συμβιβάζονται: η διαφορά ανά
στοίχημα είναι μικροσκοπική, $+0{,}00059$ στον WHC, αλλά ανατοκίζεται σε 7.216 στοιχήματα,
ενώ η αβεβαιότητα μένει αρκετά πλατιά ώστε το τεστ να μην την πιστοποιεί ($p=0{,}089$).

Το ενδιαφέρον είναι \emph{τι} μαθαίνει. Ο πολλαπλασιαστής κάθεται σχεδόν παντού πάνω από το
ένα τέταρτο, και η τιμή του ακολουθεί το πόσο κερδοφόρος είναι ο στόχος. Στον PSC, που
είναι ζημιογόνος, συνεχίζει να κατεβαίνει ως το 0,1· στον WHC ανεβαίνει ως το 0,9. Δηλαδή
ανεβάζει το ρίσκο εκεί που υπάρχει περιθώριο και το μαζεύει εκεί που δεν υπάρχει, κάτι που
μια σταθερά δεν μπορεί να κάνει.

Παράπλευρο αλλά χρήσιμο: το \emph{πλήρες} Kelly ισοπαλεί σε 16 από τα 20 κελιά και χάνει
μόνο σε 3. Το σταθερό ένα τέταρτο είναι λοιπόν πιο συντηρητικό απ' όσο χρειάζεται στους
κερδοφόρους στόχους, και ο κίνδυνος του πλήρους συγκεντρώνεται σχεδόν όλος στον PSC, όπου
μηδενίζει το κεφάλαιο. Ο μαθαινόμενος πολλαπλασιαστής είναι ο ενδιάμεσος δρόμος: κρατάει το
μεγαλύτερο μέρος της ανόδου και αποφεύγει εκείνη την καταστροφή.

\paragraph{Πού υπάρχει άκρη.} Το ποιοι bookmakers είναι καν κερδίσιμοι το χαρτογραφήσαμε σε
δέκα στόχους με το σταθερό ένα τέταρτο, ώστε οι γραμμές να μένουν συγκρίσιμες μεταξύ τους.

\begin{table}[H]
\centering
\small
\begin{tabular}{@{}llrrrl@{}}
\toprule
\textbf{Εναντίον} & \textbf{Φάση} & \textbf{Στοιχήματα} & \textbf{Κεφάλαιο} & \textbf{$p$} & \\
\midrule
WHC     & closing & 7.216  & $84{,}31\times$ & $<0{,}001$ & κέρδος \\
\midrule
BWC     & closing & 6.868  & $2{,}82\times$  & $0{,}071$  & --- \\
VC\_BVC & closing & 4.735  & $2{,}01\times$  & $0{,}185$  & --- \\
WH      & opening & 5.942  & $1{,}31\times$  & $0{,}566$  & --- \\
B365C   & closing & 7.361  & $1{,}26\times$  & $0{,}691$  & --- \\
BW      & opening & 7.666  & $1{,}03\times$  & $0{,}959$  & --- \\
B365    & opening & 10.872 & $0{,}49\times$  & $0{,}191$  & --- \\
VC\_BV  & opening & 8.284  & $0{,}44\times$  & $0{,}075$  & --- \\
PSC     & closing & 20.198 & $0{,}38\times$  & $0{,}155$  & --- \\
\midrule
PS      & opening & 27.214 & $0{,}22\times$  & $0{,}041$  & ζημιά \\
\bottomrule
\end{tabular}
\caption{Kelly $\tfrac14$, ξεκινώντας με κεφάλαιο 1, βαθμονομημένες πιθανότητες, 10 σεζόν.
Συγκρίνουμε στο $p$ του μέσου λογαριθμικού ρυθμού ανά στοίχημα. Τα τελικά κεφάλαια δεν
συγκρίνονται μεταξύ τους, γιατί ο αριθμός στοιχημάτων διαφέρει κατά τάξη μεγέθους.}
\end{table}

Σε 10 bookmakers βγάζουμε σημαντικό κέρδος σε έναν και σημαντική ζημιά σε έναν. Ο PSC
βγαίνει ισοπαλία ($p=0{,}155$), που συμφωνεί με τη σύγκριση πρόβλεψης. Η αντιστοίχιση φάσης
μετράει: με το πλήρες panel αντί για το closing-only, χειροτερεύουμε σε κάθε έναν από τους
closing στόχους και ο PSC γίνεται σημαντική ζημιά. Ο μαθαινόμενος πολλαπλασιαστής έχει
τρέξει στους πέντε closing στόχους· στους opening μένει να γίνει.

Το μήνυμα είναι ότι καλύτερο log-loss δεν μεταφράζεται σε κέρδος. Η γκανιότα είναι
μεγαλύτερη από τη διαφορά ποιότητας που κερδίσαμε, και ό,τι άκρη υπάρχει αφορά
\emph{συγκεκριμένους} bookmakers. Το μέγεθος του στοιχήματος αλλάζει πόσο γρήγορα
μεγαλώνει το κεφάλαιο εκεί που η άκρη υπάρχει· δεν δημιουργεί άκρη εκεί που δεν υπάρχει.

% =====================================================================
\section{Συμπεράσματα}
% =====================================================================

\begin{itemize}
  \item Εκπαιδευμένο ως το 2024 και τρέχοντας συνεχώς, το μείγμα προβλέπει το 2025--26
        σημαντικά καλύτερα από \emph{κάθε} τιμή ανοίγματος της αγοράς, και είναι αδιάκριτο
        από τις τιμές κλεισίματος. Δεν τις ξεπερνά.
  \item Στην πλήρη δεκαετία ο PSC μένει σημαντικά μπροστά, αλλά στη διαμόρφωση που κρατάμε
        η διαφορά παύει να είναι σημαντική.
  \item Ο OGD είναι ο καλύτερος από τους τρεις, σε κάθε panel και κανονικοποίηση. Τα Hedge
        και FTRL χάνουν από τον απλό μέσο όρο σε κάθε ομοιογενές panel.
  \item Η μάθηση αποδίδει όταν οι experts διαφέρουν σε ποιότητα. Το κέρδος δεν έρχεται τόσο
        από τον αλγόριθμο όσο από την ανομοιογένεια του panel, και αυτό είναι μάλλον το πιο
        χρήσιμο που έχουμε να πούμε σε όποιον θα εφάρμοζε το πλαίσιο αλλού.
  \item Το θεωρητικά βέλτιστο βήμα δεν βγήκε βέλτιστο στην πράξη για κανέναν από τους τρεις,
        και μάλιστα προς αντίθετες κατευθύνσεις: το OGD το ήθελε 100 φορές μεγαλύτερο, το
        Hedge τέσσερις φορές μικρότερο, το FTRL είκοσι. Τα worst-case φράγματα προστατεύουν
        από έναν αντίπαλο που οι αποδόσεις στοιχήματος δεν είναι.
  \item Καλύτερο log-loss δεν σημαίνει κέρδος. Πρέπει να ξεπεραστεί και η γκανιότα, που
        είναι μεγαλύτερη από τη διαφορά ποιότητας που κερδίσαμε.
  \item Το μέγεθος του στοιχήματος αξίζει να μαθαίνεται αντί να καρφώνεται, αρκεί να
        μαθαίνεται το σωστό πράγμα. Ο μαθαινόμενος πολλαπλασιαστής ισοπαλεί στατιστικά με το
        σταθερό ένα τέταρτο, αλλά πολλαπλασιάζει το τελικό κεφάλαιο εκεί που υπάρχει
        πραγματική άκρη και μαζεύεται μόνος του εκεί που δεν υπάρχει. Η προηγούμενη
        προσπάθεια, που μάθαινε απευθείας το κλάσμα του κεφαλαίου, έχανε συστηματικά ---
        το πρόβλημα ήταν η διατύπωση, όχι η ιδέα.
  \item Ο αλγόριθμος δεν φτάνει ποτέ σε τελική απάντηση, μαθαίνει μόνο μέσα από το να 
       διαψεύδεται, και δεν υπάρχει χωρισμός θεωρίας και πράξης.
\end{itemize}



\vspace{1em}
\hrule
\vspace{0.6em}
{\small
\textbf{Σημείωση για τα χρονικά παράθυρα.} Τα νούμερα του κυρίως κειμένου είναι στις 10
σεζόν, εκτός από την ενότητα «Πρόβλεψη για το 2025--26», που στηρίζεται σε πραγματική
χρονική τομή στις 31/12/2024.
}

\end{document}
